\chapter{Dimensional Validation and Multi-Scale Limits of Unified Forces}
\label{app:unifiedforces_validation}
\addtocontents{toc}{\protect\setcounter{tocdepth}{0}}

\textbf{Objective:} This appendix validates the Tension Field Theory (TFT) unified force equations by proving their dimensional consistency and evaluating them across the three critical geometric scales of the universe. We demonstrate that projecting the single foundational tension equation ($\tau = \rho(r)c_E^2$) successfully recovers the absolute theoretical limits of the vacuum and the experimental coupling constants of observable matter.

\section{The Master Geometric Law \& Torsional Invariant}
The density at any given scale is strictly governed by the geometric dilution of the Substrate, a necessity to mathematically conserve the universal Torsional Invariant ($J_i = \rho_i r_i^4$) across all scales:
\begin{equation}
    \rho(r) = \mathcal{F} \cdot \rho_S \cdot \left(\frac{\ell_S}{r}\right)^4
\end{equation}

\textbf{Where:}
\begin{itemize}
    \item $\rho_S$: Absolute Substrate density ($\approx 5.1549 \times 10^{96} \text{ kg/m}^3$).
    \item $\ell_S$: Fundamental Substrate Lattice Spacing ($\approx 1.6163 \times 10^{-35} \text{ m}$).
    \item $\mathcal{F}$: Topological Form Factor ($1.5$ for the Soft Floor active vacuum; $1.0$ for fundamental nucleons).
\end{itemize}

\textbf{Additional Geometric Projection Parameter:}
\begin{itemize}
    \item $\kappa$: Volumetric Configuration Parameter ($1$ for Simple Cubic grid; $4\pi/3$ for spherical topologies).
\end{itemize}

\section{Dimensional Validation of the Master Equations}
Before evaluating the equations at specific scales, we must verify that the geometric projections of the foundational tension equation strictly reproduce the correct dimensional units. Because the unified framework projects a single tension ($\tau$) through different geometric dimensions (1D, 2D, 3D), the resulting phenomenological constants naturally possess distinct dimensional footprints.

\begin{itemize}
    \item \textbf{Base Tension ($\tau$):}
\end{itemize}
\begin{equation}
    \tau = \rho c_E^2 \to [\text{kg} \cdot \text{m}^{-3}][\text{m}^2 \cdot \text{s}^{-2}] = [\text{kg} \cdot \text{m}^{-1} \cdot \text{s}^{-2}] = [\text{N} \cdot \text{m}^{-2}] \quad (\text{Stress / Pressure})
\end{equation}

\begin{itemize}
    \item \textbf{Strong Force ($\sigma_s$ / 1D Bond Pull):} Projecting tension through a cross-sectional area ($A_{chan}$):
\end{itemize}
\begin{equation}
    \sigma_s = \tau A_{chan} \to [\text{N} \cdot \text{m}^{-2}][\text{m}^2] = [\text{N}] \quad (\text{Classical Force})
\end{equation}

\begin{itemize}
    \item \textbf{Electromagnetism ($\alpha$ / 2D Surface Twist):} Evaluating the Absolute Lattice Shear Coupling ($K_L$, measured in $\text{J}\cdot\text{m}$) against the spatial projection of the action limit ($\hbar c_E$):
\end{itemize}
\begin{equation}
    \alpha = \frac{K_L}{\hbar c_E} \to \frac{[\text{J} \cdot \text{m}]}{[\text{J} \cdot \text{s}][\text{m} \cdot \text{s}^{-1}]} = \frac{[\text{J} \cdot \text{m}]}{[\text{J} \cdot \text{m}]} = 1 \quad (\text{Dimensionless})
\end{equation}

\begin{itemize}
    \item \textbf{Weak Force ($G_F$ / 3D Gear Slip):} Projecting tension through a defect reconfiguration volume ($V$):
\end{itemize}
\begin{equation}
    G_F = (\tau V)^{-2} \to \left([\text{N} \cdot \text{m}^{-2}][\text{m}^3]\right)^{-2} = [\text{N} \cdot \text{m}]^{-2} = [\text{Energy}]^{-2}
\end{equation}
\textit{(Note: This derives the operational high-energy form of the Fermi constant ($\text{GeV}^{-2}$) utilized in quantum field theory, representing a pure energy-yield threshold).}

\begin{itemize}
    \item \textbf{Gravity ($\alpha_G$ / Restorative Radial Tension):} Evaluating the dimensionless gravitational coupling ratio:
\end{itemize}
\begin{equation}
    \alpha_G = \frac{\kappa \rho G r^2}{c_E^2} \to \frac{[\text{kg} \cdot \text{m}^{-3}][\text{m}^3 \cdot \text{kg}^{-1} \cdot \text{s}^{-2}][\text{m}^2]}{[\text{m}^2 \cdot \text{s}^{-2}]} = \frac{[\text{s}^{-2}]}{[\text{s}^{-2}]} = 1 \quad (\text{Dimensionless})
\end{equation}

\section{Test 1: The Hard Floor Limit (State 0: The Latent Substrate)}
At the absolute structural limit of the universe, the radius is the Substrate Lattice Spacing ($r = \ell_S$) and the density is the unbroken Substrate density ($\rho_S = c_E^2 / G\ell_S^2$). Because the Latent Substrate is a perfectly rigid Simple Cubic (SC) grid, topological wave geometry ceases to exist. All structural projections revert strictly to rectilinear bounds: $\kappa = 1, A_{chan} = \ell_S^2, V = \ell_S^3$.

\vspace{0.2cm}
\noindent \textbf{Strong Force ($\sigma_s$):} 1D Confinement Flux across a single SC lattice face ($A_{chan} = \ell_S^2$):
\begin{equation}
    \sigma_s = \rho_S c_E^2 \ell_S^2 = \left(\frac{c_E^2}{G\ell_S^2}\right) c_E^2 \ell_S^2 = \frac{c_E^4}{G}
\end{equation}
\noindent \textbf{Result:} Strictly evaluates to the ultimate theoretical yield tension of the Hard Floor ($F_S$).

\vspace{0.2cm}
\noindent \textbf{Electromagnetism ($\alpha$):} The ultimate substrate capacity limit. Evaluating the absolute spatial tension ($K_L$) against the Substrate temporal limit. As established in \hyperref[app:geometric_alpha]{Appendix \ref*{app:geometric_alpha}}, in the completely rigid Substrate, transverse shear tension ($K_L$) and temporal action ($\hbar c_E$) are perfectly merged and equivalent potentials.
\begin{equation}
    \alpha = \frac{K_L}{\hbar c_E} = 1
\end{equation}
\noindent \textbf{Result:} The math physically proves that the topological gear ratio locks at the absolute structural limit, extinguishing independent Charge and Light. The electromagnetic coupling is strictly $1$.

\vspace{0.2cm}
\noindent \textbf{Weak Force ($G_F$):} 3D Defect Reconfiguration across a fundamental cubic cell ($V = \ell_S^3$):
\begin{equation}
    G_F = (\rho_S c_E^2 \ell_S^3)^{-2} = \left[ \left(\frac{c_E^2}{G\ell_S^2}\right) c_E^2 \ell_S^3 \right]^{-2} = (F_S \ell_S)^{-2}
\end{equation}
\noindent \textbf{Result:} Perfectly resolves to the exact inverse square of the fundamental substrate energy quantum ($E_S^{-2}$).

\vspace{0.2cm}
\noindent \textbf{Gravity ($\alpha_G$):} Macroscopic anchoring saturation.
\begin{equation}
    \alpha_G = \frac{(1) \rho_S G \ell_S^2}{c_E^2} = \frac{\left(\frac{c_E^2}{G\ell_S^2}\right) G \ell_S^2}{c_E^2} = 1
\end{equation}
\noindent \textbf{Result:} Mathematically proves Postulate I. Gravity bottoms out exactly at absolute structural saturation ($\alpha_G = 1$).

\section{Test 2: The Soft Floor Limit (State I: The Active Vacuum)}
At the scale of the dynamic active vacuum, the radius is the FCC Node Radius ($r = r_D$). The physical lattice density ($\rho_D$) utilizes a spherical form factor ($\mathcal{F} = 1.5$) applied to the quantum energy limit. Because the Soft Floor represents topological standing waves, projection channels adopt circular/spherical geometry ($\kappa = 4\pi/3, A_{chan} = 2\pi r_D^2, V = 2\pi r_D^3$).

\vspace{0.2cm}
\noindent \textbf{Strong Force ($\sigma_s$):} Evaluating the 1D tension across a cylindrical flux tube ($A_{chan} = 2\pi r_D^2$) using the spherical lattice density ($\rho_D = 1.5 \cdot \hbar / c_E r_D^4$):
\begin{equation}
    \sigma_s = \rho_D c_E^2 (2\pi r_D^2) = \left(\frac{1.5\hbar}{c_E r_D^4}\right) c_E^2 (2\pi r_D^2) = \frac{3\pi\hbar c_E}{r_D^2}
\end{equation}
\noindent \textbf{Result:} Because the fundamental quantum energy limit is $E_Q = \hbar c_E / r_D$, the confinement force on the Soft Floor resolves identically to the fundamental quantum energy gradient: $\sigma_s = 3\pi(E_Q/r_D)$.

\vspace{0.2cm}
\noindent \textbf{Electromagnetism ($\alpha$):} As derived geometrically in \hyperref[app:geometric_alpha]{Appendix \ref*{app:geometric_alpha}}, the active vacuum creates a topological boundary-loss limit resulting in the fine-structure constant ($\alpha \approx 1/137.036$). Evaluating the derived Absolute Lattice Shear Coupling ($K_L \approx 2.3071 \times 10^{-28} \text{ J}\cdot\text{m}$) against the quantum action limit ($\hbar c_E \approx 3.1615 \times 10^{-26} \text{ J}\cdot\text{m}$):
\begin{equation}
    \alpha = \frac{K_L}{\hbar c_E} \approx \frac{2.3071 \times 10^{-28}}{3.1615 \times 10^{-26}} \approx \frac{1}{137.036}
\end{equation}
\noindent \textbf{Result:} Perfectly validates the geometric derivation. The spatial tension of the Soft Floor scales strictly by the 144-channel topological gear ratio.

\vspace{0.2cm}
\noindent \textbf{Weak Force ($G_F$):} Evaluating the 3D defect reconfiguration across the topological defect volume ($V = 2\pi r_D^3$):
\begin{equation}
    G_F = (\rho_D c_E^2 \cdot 2\pi r_D^3)^{-2} = \left[ \left(\frac{1.5\hbar}{c_E r_D^4}\right) c_E^2 (2\pi r_D^3) \right]^{-2} = \left(\frac{3\pi\hbar c_E}{r_D}\right)^{-2}
\end{equation}
\noindent \textbf{Result:} The Weak Force perfectly resolves to the exact inverse square of the dynamic quantum energy limit: $G_F = (3\pi E_Q)^{-2}$.

\vspace{0.2cm}
\noindent \textbf{Gravity ($\alpha_G$):}
\begin{equation}
    \alpha_G = \frac{\left(\frac{4\pi}{3}\right) \cdot \rho_D G r_D^2}{c_E^2} = \frac{\left(\frac{4\pi}{3}\right) \cdot 1.5 \left(\frac{\hbar}{c_E r_D^4}\right) G r_D^2}{c_E^2}
\end{equation}
Because $(4/3) \cdot 1.5 = 2$, the equation simplifies flawlessly:
\begin{equation}
    \alpha_G = \frac{2\pi\hbar G}{c_E^3 r_D^2} = 2\pi \left(\frac{\ell_S}{r_D}\right)^2
\end{equation}
\noindent \textbf{Result:} The emergence of the $2\pi$ multiplier perfectly mirrors the transition into wave topology.

\section{Test 3: The Nucleon Limit (Composite Matter)}
To evaluate the forces at the macroscopic scale, we calculate the effective density ($\rho_{eff}$) of a fundamental nucleon ($\mathcal{F}=1.0$) at the mechanical nucleon radius ($r_N \approx 0.88 \times 10^{-15} \text{ m}$):
\begin{equation}
    \rho_{eff} = 1.0 \cdot (5.1549 \times 10^{96}) \cdot \left(\frac{1.6163 \times 10^{-35}}{0.88 \times 10^{-15}}\right)^4 \approx 5.8665 \times 10^{17} \text{ kg/m}^3
\end{equation}

\vspace{0.2cm}
\noindent \textbf{Strong Force ($\sigma_s$):}
\begin{equation}
    \sigma_s = \rho_{eff} c_E^2 (2\pi r_N^2) = (5.8665 \times 10^{17})(8.9876 \times 10^{16}) \cdot 2\pi (0.88 \times 10^{-15})^2 \approx 2.5655 \times 10^5 \text{ N}
\end{equation}
\noindent \textbf{Result:} $2.5655 \times 10^5 \text{ N} \approx 1.6 \text{ GeV/fm}$, perfectly recovering the experimental QCD string tension (confinement force) limits of $1.0 - 2.0 \text{ GeV/fm}$.

\vspace{0.2cm}
\noindent \textbf{Electromagnetism ($\alpha$):} Because the Nucleon is a vortex loop, its mass-moment governs the localized action ($\hbar = c_E J_N$). The Nucleon Torsional Invariant is $J_N = \rho_{eff} r_N^4 \approx 3.5181 \times 10^{-43} \text{ kg}\cdot\text{m}$. Validating the absolute spatial tension against this localized anchor:
\begin{equation}
    \alpha = \frac{K_L}{c_E^2 J_N} = \frac{2.3071 \times 10^{-28}}{(8.9876 \times 10^{16}) \cdot (3.5181 \times 10^{-43})} \approx \frac{1}{137.05}
\end{equation}
\noindent \textbf{Result:} This is the mathematical proof of scale invariance. The localized geometric slack perfectly dictates the macroscopic coupling strength, evaluating strictly to $\approx 1/137$ at the atomic level.

\vspace{0.2cm}
\noindent \textbf{Weak Force ($G_F$):} Evaluating the 3D reconfiguration at the heavily focused weak interaction vertex ($r_w \approx 5 \times 10^{-18} \text{ m}$). By the master scaling law ($\rho \propto r^{-4}$), the localized density spikes exponentially at this compressed scale to $\rho(r_w) \approx 5.6290 \times 10^{26} \text{ kg/m}^3$:
\begin{equation}
    E_{weak} = 2\pi \cdot \rho(r_w) c_E^2 r_w^3 \approx 3.9733 \times 10^{-8} \text{ J} \approx 248 \text{ GeV}
\end{equation}
\begin{equation}
    G_F = (E_{weak})^{-2} \approx (248 \text{ GeV})^{-2} \approx 1.6 \times 10^{-5} \text{ GeV}^{-2}
\end{equation}
\noindent \textbf{Result:} Correctly recovers the exact Weak Interaction Vacuum Expectation Value ($\approx 246 \text{ GeV}$) and the correct standard magnitude of the Fermi constant ($\sim 10^{-5} \text{ GeV}^{-2}$).

\vspace{0.2cm}
\noindent \textbf{Gravity ($\alpha_G$):} Using the spherical configuration ($\kappa = 4\pi/3 \approx 4.1888$):
\begin{equation}
    \alpha_G = \frac{(4.1888)(5.8665 \times 10^{17})(6.6743 \times 10^{-11})(0.88 \times 10^{-15})^2}{(299792458)^2} \approx 1.4132 \times 10^{-39}
\end{equation}
\noindent \textbf{Result:} Flawlessly recovers the magnitude of the standard Newtonian gravitational coupling constant.